\section{Discussion}

The operating termination determines the component freedom available for cancellation. A fixed $N$-stage reference branch admits an $N$-dimensional open-circuit matching family and an $(N-1)$-dimensional family at one positive resistive load. Two distinct known terminations fix the complete polynomial pair and therefore the components. The rational-load result further separates the effect of circuit order from the effect of load type.

The coefficient representation gives a constructive design sequence. Specify the source and output measurement planes and the terminal admittance. Select the required transfer coefficients, choose the remaining free coefficients, and apply the recursive positive-realization test in Eq.~\eqref{eq:general-inverse}. Evaluate component and load perturbations on the recovered circuit. The P/Q and P/T constructions illustrate why an open-circuit equivalence or a single-load match must be tested at the intended operating termination.

\citet{casas1996basics} obtained cancellation through floating coupling in a different topology, including higher-order extensions. The grounded ladders here retain shunt current paths, and the pair $(A_N,B_N)$ exposes terminal loading. \citet{hornero2021cm} analyzed frequency-dependent mismatch through the common-mode response. The present representation connects that response to a positive component family and then to the calibrated ratio observable.

Strict component positivity preserves the stated order and the uniqueness result. For example, setting $q_1=0$ in the second-order ladder merges two series resistors. $(1,2,0,3)$ and $(3/2,3/2,0,3)$ yield the same coefficients $(9,0,3,0)$. The inverse then loses the separation of the two resistances. Unknown source resistance similarly combines with the first series resistor unless the source contribution is measured independently.

\subsection{Dynamic fidelity and calibration}

Direct-current calibration changes which transfer functions need to match. A constant-resistance ratio rejects all admissible current histories when the normalized channel responses are equal. Raw common-mode cancellation additionally requires equal absolute gains. The matching dimensions in Table~\ref{tab:load-dimensions} concern the raw response and retain that gain requirement.

For changing resistance, matched channels preserve a weighted history average. Reducing the direct-current group delay $\tau=D'(0)/D(0)$ at the operating load reduces the first-moment tracking bound, while reducing the channel-delay mismatch $\Delta\tau=\tau_P-\tau_N$ (Section~\ref{sec:dynamic-mismatch}) suppresses the additional first-order current-slew term. Tracking delay and channel matching are different design objectives. The three- and four-stage transient examples quantify both effects. Identical channels retained current-dependent history weighting, and capacitor mismatch generated the predicted signed error during excitation changes.

The two-load result supplies an identification strategy for $2N$ component parameters. In an experiment, complex responses under known loads would determine the polynomial pair with uncertainty intervals. Closely spaced loads and strongly attenuated responses can amplify errors in that recovery, as described in Section~\ref{sec:general-order}. Measurement-band selection and conditioning therefore accompany the exact identifiability result.

\subsection{Scope of the evidence and experimental limitations}

The evidence in this paper consists of exact algebra and independent circuit simulation; no physical measurements are included. A bench validation of the termination-dependent families would require measured component values, characterized complex source and probe impedances, separate branch responses under known terminations, common-mode and differential excitation of the pair, repeatability records, and an uncertainty budget sufficient to separate nominal from measured-component predictions.

Dynamic validation additionally requires the two source signals and both filtered outputs on a common time base, with constant-resistance and changing-resistance criteria evaluated separately and the measured waveforms driving the reference model, so that excitation shape and channel mismatch remain distinguishable. Leakage measurements near the instrument floor are bounded by the corresponding detection limit rather than resolved.

